\subsection{局部凸空间上的前测度}

设 E 是局部凸空间。记 \(\mathcal{F}(\mathrm{E})\) 为 E 的有限余维闭线性子空间所成的集合，并按关系 \(\supset\) 排序。对每个 \(V \in \mathcal{F}(\mathrm{E})\)，\(p_{V}\) 表示从 E 到 E/V 的典范映射。设 V 和 W 是 \(\mathcal{F}(\mathrm{E})\) 的两个元素，满足 \(V \supset W\)；记 \(p_{VW}\) 为从 E/W 到 E/V 的映射，它由 E 的恒等映射通过取商诱导。族 \(\mathcal{Q}(\mathrm{E}) = (\mathrm{E}/\mathrm{V}, p_{\mathrm{VW}})\) 是以 \(\mathcal{F}(\mathrm{E})\) 为指标的局部凸空间逆系统。称其为 E 的有限维商空间逆系统。

可以证明，逆系统 \(\mathcal{Q}(\mathbf{E})\) 的逆极限典范同构于 \(E^{\prime*}\)，即 \(E^{\prime}\) 的代数对偶，并赋予弱拓扑 \(\sigma(\mathbf{E}^{\prime*},\mathbf{E}^{\prime})\)。

定义 1. --- 设 E 是局部凸空间。称 E 上的前测度为有限维商空间逆系统上的每个测度逆系统 \(^{(1)}\)（§4, No.~2, 定义 1）。

换言之，前测度 \(\mu\) 定义在 E 上，是族 \((\mu_{\mathrm{V}})_{\mathrm{V}\in \mathcal{F}(\mathrm{E})}\)，其中 \(\mu_{\mathrm{V}}\) 是有限维空间 \(\mathrm{E / V}\) 上的有界（正）测度，并且 \(\mu_{\mathrm{V}} = p_{\mathrm{VW}}(\mu_{\mathrm{W}})\) 当 \(\mathrm{V}\supset \mathrm{W}\) 时成立。所有测度 \(\mu_{\mathrm{V}}\) 都有相同的总质量，该质量称为前测度 \(\mu\) 的总质量。

要使 E 的子空间 V 属于 \(\mathcal{F}(\mathrm{E})\)，必要且充分的条件是：存在有限个元素 \(x_{1}^{\prime},\ldots,x_{n}^{\prime}\)（属于 \(E^{\prime}\)），使得 V 由满足 \(x\in E\) 且 \(\langle x,x_{i}^{\prime}\rangle=0\) 对 \(1\leqslant i\leqslant n\) 成立的 x 所成（TVS, II, §6, No.~3, Th. 1 的推论 2 以及 No.~5, 命题 7 的推论 2）。此外，有限维向量空间上存在唯一的豪斯多夫拓扑向量空间拓扑（TVS, I, §2, No.~3, 定理 2）。因此，E 上前测度的概念只依赖于 E 的对偶 \(E^{\prime}\)。

设 \(\lambda\) 是 E 上的有界测度。对每个 \(V \in \mathcal{F}(E)\)，记 \(\widetilde{\lambda}_{V}\) 为 \(\lambda\) 在典范映射 \(p_{V}\) 下的像。由 \(p_{V} = p_{\mathrm{vw}} \circ p_{W}\)，对于 \(\mathcal{F}(E)\) 中满足 \(V \supset W\) 的任意两个元素 V 和 W；因此，族 \(\widetilde{\lambda} = (\widetilde{\lambda}_{V})_{V \in \mathcal{F}(E)}\) 是 E 上的前测度。称 \(\widetilde{\lambda}\) 为与测度 \(\lambda\) 相对应的前测度。立即可见，\(\lambda\) 与 \(\widetilde{\lambda}\) 的总质量相同。

命题 1. --- 设 E 是局部凸空间。映射 \(\lambda \mapsto \widetilde{\lambda}\) 将 E 上的有界测度集合双射到满足以下条件的 E 上的前测度 \((\mu_{\mathrm{V}})_{\mathrm{V} \in \mathcal{F}(\mathrm{E})}\) 所成的集合：

对每个 \(\varepsilon > 0\)，存在紧子集 \(\mathbf{K}\)（属于 \(\mathbf{E}\)），使得 \(\mu_{\mathrm{V}}\big(\mathrm{E} / \mathrm{V} - p_{\mathrm{V}}(\mathrm{K})\big) \leqslant \varepsilon\) 对所有 \(\mathrm{V} \in \mathcal{F}(\mathrm{E})\) 成立。

已知 E 上连续线性形式的核的交等于 0（TVS, II, §4, No.~1, 命题 2 的推论 1）；因此 \(\bigcap_{\mathrm{V} \in \mathcal{F}(\mathrm{E})} \mathrm{V} = \{0\}\)，且族 \((p_{\mathrm{V}})_{\mathrm{V} \in \mathcal{F}(\mathrm{E})}\) 相容且分离。该

命题于是由 §4 第 2 号定理 1 得出。

特别地，映射 \(\lambda \mapsto \widetilde{\lambda}\) 是单射。若 \(\mu\) 是 \(\mathbf{E}\) 上的前测度，且 \(\lambda\) 上存在有界测度 \(\mathbf{E}\) 使得 \(\mu = \widetilde{\lambda}\)，则我们滥用语言称 \(\mu\) 为测度。若 \(\mathbf{E}\) 是有限维的，则每个前测度 \(\mu = (\mu_{\mathrm{V}})_{\mathrm{V} \in \mathcal{F}(\mathbf{E})}\) 都是测度：因为 \(\{0\} \in \mathcal{F}(\mathbf{E})\)、\(\mathrm{E}/\{0\} = \mathrm{E}\) 且 \(p_{\mathrm{V},\{0\}} = p_{\mathrm{V}}\)，从而有 \(\mu_{\mathrm{V}} = p_{\mathrm{V}}(\mu_{\{0\}})\) 对所有 \(\mathrm{V} \in \mathcal{F}(\mathrm{E})\) 成立；换言之，\(\mu = \widetilde{\lambda}\)，其中 \(\lambda = \mu_{\{0\}}\)。

命题 2. --- 设 T 是可数集，E 是 T 上实函数所成的向量空间，赋予逐点收敛拓扑。E 上的每个前测度都是测度。

对每个 \(t \in \mathbf{T}\)，令 \(\varepsilon_t\) 为 \(f \mapsto f(t)\) 上的线性形式 \(\mathbf{E}\)。已知（TVS, II, §6, No.~6, 命题 8 的推论 2），族 \((\varepsilon_t)_{t \in \mathbf{T}}\) 是向量空间 \(\mathbf{E}'\) 的一个基。记 \(\Phi\) 为 ` I' 的有限子集所成的集合，对每个 \(J \in \Phi\)，令 \(\mathbf{E}_J\) 为 \(\mathbf{T}\) 上在 \(J\) 的每一点都为零的函数所成的集合。设 \(F \in \mathcal{F}(\mathbf{E})\)；由于 \(F^\circ\) 是 \(F\) 的正交补，是 \(E'\) 的有限维子空间，存在 \(J \in \Phi\)，使得 \(F^\circ\) 包含于线性子空间 \(G\) 中；该子空间是 \(E'\) 中由 \(\varepsilon_t\) 对 \(t \in J\) 生成的。由于 \(F^\circ \subset G\)，有 \(E_J = G^\circ \subset F^{\circ\circ} = F\)，且可数族 \((E_J)_{J \in \Phi}\) 在 \(\mathcal{F}(\mathbf{E})\) 中共尾。命题于是由 §4 第 3 号定理 2 得出。

